\subsection{The explicit dimension curve}\label{hardgap:dimension}
Put
\[
 \alpha=\frac{3-\sqrt7}{4},\qquad
 a_c=\frac{\sqrt6-2}{4},\qquad
 b_c(a)=\frac{1+2a}{4+2a},\qquad
 b_F(a)=\frac{2a}{1-2a}.
\]
Thus \(d_\alpha=1+\alpha\), \(d_c=1+a_c\). The relevant threshold
can now be evaluated exactly, rather than by compact optimization.

First consider \(\alpha\le a\le a_c\) and \(2a\le b\le b_c(a)\).
The two-gap guard holds strictly throughout this region. Indeed,
\(\partial_b R=8(1+a)b-4(a+2)<0\) for \(b\le1/2\), and
\[
 R(a,2a)=16a^3+8a^2-13a+1.
\]
This last polynomial decreases for \(\alpha\le a\le1/6\), while
\[
 R(\alpha,2\alpha)=3(11\alpha-1)<0.
\]
The last sign follows from \(\alpha<1/11\), equivalently
\(29<11\sqrt7\), whose squares are \(841<847\).

For fixed \(a\), the derivative of \(C_2\) on this region is
\[
 \partial_bC_2=
 \frac{2+11a+8a^2-(8+14a+8a^2)b}{(1-a)(1+2b)^3}>0.
\]
Its numerator decreases with \(b\), and multiplying its value at
\(b_c(a)\) by \(4+2a\) gives \(18a(1+a)>0\).
At the lower endpoint,
\begin{equation}\label{hardgap:baseline-sign}
 C_2(a,2a)-a
 =\frac{a(1+2a)(8a^2-12a+1)}{(1-a)(1+4a)^2}.
\end{equation}
This vanishes at \(a=\alpha\) and is negative for
\(\alpha<a\le a_c\). At the upper endpoint, \(C_2=C_1\), and
\[
 C_1(a,b)<a\quad\Longleftrightarrow\quad b<b_F(a).
\]
The equality \(b_F(a)=b_c(a)\) occurs at \(a=a_c\), since it is
equivalent to \(8a^2+8a-1=0\). For \(a<a_c\), the left side is
smaller; consequently \(C_2(a,b_c(a))>a\).

It follows that \(C_2(a,b)=a\) has a unique root in
\([2a,b_c(a)]\), with the appropriate endpoint equalities. Clearing
positive denominators gives
\[
 (4a^2-8a-2)b^2+(8a^2-a+2)b-3a=0.
\]
Its smaller root is exactly \(r_2(a)\) in
\eqref{hardgap:rootformula}. Existence of the crossing and the strictly
positive derivative at it show that the discriminant is positive;
the stated radical is therefore real on the full closed interval.
The root is continuous and strictly larger than \(2a\) for
\(\alpha<a\le a_c\). Its endpoint values are
\[
 r_2(\alpha)=2\alpha,\qquad
 r_2(a_c)=b_c(a_c)=\frac{\sqrt6-1}{5}.
\]
Thus the dimension branches agree at the two joins; at the second join
their common value is \((4+\sqrt6)/5\).

For \(a_c\le a\le1/4\), the cutoff is \(b_F(a)\). To verify that
no lower-slope interval is missed, suppose \(2a\le b\le b_c(a)\).
The inequality \(2a\le b_c(a)\) forces
\(a\le(\sqrt{13}-3)/4<1/6\). The same guard calculation above applies,
and monotonicity gives
\[
 C_2(a,b)\le C_2(a,b_c(a))=C_1(a,b_c(a))\le a,
\]
strictly except at \(a=a_c,b=b_c(a_c)\).
For \(b\ge b_c(a)\), the single-collapse upper bound gives
\(C_1<a\) precisely for \(b<b_F\); above \(b=1/2\) the earlier
high-slope estimate supplies the same bound. The cases \(b<2a\)
are already covered by the coherent criterion.

\begin{proof}[Proof of Theorem~\ref{target:hardgap}]
For \(1<d\le d_\alpha\), or whenever \(D<2d-1\), use
Proposition~\ref{prop:coherentborel}. In every remaining case of
\(D<B_{\rm H}(d)\), the preceding calculation gives a universal
profile cost strictly below \(a=d-1\), using \(C_2\) or \(C_1\)
with \(b=D-1\). The two formulas agree at \(b=b_c(a)\), and the
two-gap guard is strict wherever that branch is used here.
Continuity and strictness therefore allow \(s<d\), \(u>D\), with
\(1<s<5/4\), \(s<u<2\), still in the certified region and with
the corresponding cost below \(s-1\).

The hard-gap profile theorem gives the same bounded number of edges,
admissibility and \(O(KT+K\epsilon N)\) perturbation error required
by Theorem~\ref{fprof:main}. Substitute its cost in that proof.
The compact restrictions, finite regularization, deletion, Fourier
iteration and summable shell limit are unchanged. They give a compact
source subset and a family of pins in \(E\) of full measure for the
selected Frostman pin probability, with absolutely continuous pinned
distance measures. In particular one
\(y\in E\) has \(|\Delta_y(E)|>0\).
\end{proof}

\subsection{Precisely what is optimal here}
For \(0<a\le1/4\), \(a\le b\le1\), define
\[
 \mathcal M(a,b)=\sup_g\lim_{\zeta\downarrow0}
 \inf_{\substack{\text{finite admissible chains}\\1-\zeta\ \text{to }0}}
       \sum c_g,
\]
where the supremum is over all 1-Lipschitz profiles with
\(ax\le g(x)\le bx\). Consider the following precisely specified
combination of sufficient tests:
\begin{equation}\label{hardgap:model-criterion}
 b<2a\qquad\text{or}\qquad\mathcal M(a,b)<a.
\end{equation}
The first is the coherent test; the second is the limiting universal
chain-cost test used in the analytic transfer. The new curve is the
exact threshold for this combination:
\begin{equation}\label{hardgap:model-optimality}
 \eqref{hardgap:model-criterion}
 \quad\Longleftrightarrow\quad
 1+b<B_{\rm H}(1+a).
\end{equation}

We have proved the sufficient direction. For the converse, when
\(0<a\le\alpha\), the two-collapse profile with upper slope \(2a\)
has cost at least \(a\), by \eqref{hardgap:baseline-sign}.
It is also admissible under every larger upper barrier, so
\(\mathcal M(a,b)\ge a\) for \(b\ge2a\).
When \(\alpha<a<a_c\), the matching two-collapse profiles give
\(\mathcal M\ge C_2\ge a\) from \(r_2(a)\) to \(b_c(a)\).
Above \(b_c(a)\), the single-collapse profiles give
\(\mathcal M\ge C_1\ge a\). Finally, for \(a\ge a_c\), those
same single-collapse profiles give \(\mathcal M\ge a\) whenever
\(b\ge b_F(a)\). The profile lower bounds include equality at all
displayed thresholds. This proves \eqref{hardgap:model-optimality}.

This optimality concerns exactly \eqref{hardgap:model-criterion}.
It gives no set \(E\) whose pinned distances have zero length.
Better analytic information, a more selective family of profiles, or a
different distance-set argument could still enlarge the sufficient region.
In particular, it does not determine the weakest possible Hausdorff--packing
condition for the distance theorem itself.

For comparison with the preceding manuscript, its profile constructions
give universal costs, so \(\mathcal M\le C_{\rm A}\) wherever that
function was used. Equation~\eqref{hardgap:model-optimality} therefore
shows that \(B_{\rm H}\) includes the entire preceding curve
\(B_{\rm A}\). The new curve also has a direct closed-form description.

Two exact examples illustrate the strict inequalities:
\[
 C_2\left(\frac9{100},\frac{181}{1000}\right)=\frac{4619}{51529},
 \qquad\frac9{100}-\frac{4619}{51529}=\frac{1861}{5152900}>0,
\]
\[
 C_1\left(\frac3{20},\frac25\right)=\frac5{36},
 \qquad\frac3{20}-\frac5{36}=\frac1{90}>0.
\]
Also \(r_2(1/10)=5/23\), so the exact cutoff at Hausdorff dimension
\(11/10\) is \(28/23\); at dimension \(23/20\) it is \(10/7\).
